% ============================================================
%  Visual Dependency Graphs Improve LLM Cross-Module Bug Fixing
%  Preprint — Overleaf-ready
% ============================================================
\documentclass[11pt,a4paper]{article}
\usepackage[margin=2.5cm]{geometry}
\usepackage{graphicx,booktabs,hyperref,amsmath,multirow,enumitem,xcolor}
\usepackage[backend=biber,style=numeric]{biblatex}
\hypersetup{colorlinks=true,linkcolor=blue,citecolor=blue,urlcolor=blue}
\title{\bfseries Visual Dependency Graphs Improve LLM Cross-Module Bug Fixing\\[3pt]
       \large Teaching Models to ``Step Back and Look at the Canvas''}
\author{}
\date{July 2026}

\begin{document}
\maketitle

\begin{abstract}
Large language models (LLMs) exhibit a well-documented tendency toward \textit{local reasoning} when performing code modifications: they focus on the single file or line implicated in an error, ignoring cross-module dependencies that may also be affected. We hypothesize that providing a \textbf{visual dependency graph}---a node-and-edge diagram of module relationships---acts as a ``global canvas'' that nudges the model toward system-level impact analysis. We test this hypothesis through five controlled experiments, training a Qwen2-VL-7B model on 334 bug-fix scenarios across 73 real-world Python projects. Results show that a dependency graph paired with a minimal text key raises cross-module file identification from 0\% to 100\% on unseen test projects, and Chain-of-Thought fine-tuning on diagram descriptions causes a +9-point net improvement in diagram utility. We release a multi-language CLI tool that auto-generates dependency graphs and text keys for injection into any LLM conversation.
\end{abstract}

% ============================================================
\section{Introduction}

When a human painter works on a detailed section of a canvas, they periodically step back to view the whole composition---a habit that prevents local distortions from accumulating into global errors. LLMs, trained predominantly on local bug-fix examples scraped from the web, lack this instinct. Given an error traceback pointing to one file, they tend to fix only that file, ignoring adjacent modules that call the changed function or depend on the altered interface.

We call this the \textit{local-repair bias}. It is not a consequence of limited context windows---the entire codebase can be provided---but rather a learned behavioral pattern. The model has been trained to act as ``a painter who never steps back.''

This paper investigates a simple intervention: providing a \textbf{module-level dependency graph} as a visual scaffold during the bug-fix prompt. The graph uses color to distinguish the changed module (red) from its direct dependents (yellow), giving the model an immediate, parallel view of the impact surface.

Our contributions are:
\begin{enumerate}[nosep]
\item A controlled experimental framework for measuring the local-repair bias and the corrective effect of dependency graphs.
\item Evidence that Chain-of-Thought (CoT) fine-tuning on diagram descriptions internalizes the ``step back'' habit, transferable to text-only settings.
\item A practical architecture---visual graph + minimal text key---that sidesteps the OCR limitations of small vision models.
\item An open-source multi-language CLI tool (\texttt{depgraph}) for generating dependency graphs and prompt-ready context snippets.
\end{enumerate}

% ============================================================
\section{Methodology}

\subsection{Dependency Graph Generation}
We use AST-based import analysis to construct a directed graph $G=(V,E)$ where each node $v \in V$ represents a source module and each directed edge $(u,v) \in E$ indicates that module $u$ imports from module $v$. The graph is rendered as a color-coded network diagram: the \textbf{changed module} is red, its \textbf{direct dependents} are yellow, and unrelated modules are gray. A compact text key maps module names to their dependency roles.

\subsection{Training Data}
We collected 334 training samples from 73 real-world Python projects (requests, click, httpx, rich, starlette, PIL, etc.). Each sample consists of:
\begin{itemize}[nosep]
\item A \textbf{bug description}: a function signature change in one module.
\item \textbf{Source code} of the changed module.
\item A \textbf{dependency diagram} PNG with color-coded dependencies.
\item An \textbf{assistant response} showing impact analysis and fixes for all affected files.
\end{itemize}
Samples are quality-filtered: 255 perfect (100\% file coverage) and 79 good ($\geq$75\% coverage). The dataset is 172 MB in JSONL format.

\subsection{Training Configuration}
We fine-tune Qwen2-VL-7B-Instruct using QLoRA (4-bit quantization, rank 8, alpha 16) for 5 epochs on a single RTX 4090 (24 GB). Training takes approximately 35 minutes per run. CoT variants prepend a structured diagram description (``The red module is X. Yellow dependents are: A, B, C...'') to every assistant response.

\subsection{Evaluation Metrics}
For each test case, we measure \textbf{file coverage}: the fraction of ground-truth affected files that the model correctly identifies in its response. We report both raw counts and percentage improvements. Three evaluation settings are used:
\begin{enumerate}[nosep]
\item \textbf{Hard matrix}: 19 bug scenarios with no dependency hints in the text prompt. The diagram is the sole source of dependency information.
\item \textbf{Toy project test}: A custom 8-module task management system, never seen during training.
\item \textbf{Cross-model validation}: MiMo-V2.5 (vision-capable API model) tested on the same scenarios.
\end{enumerate}

% ============================================================
\section{Experiments}

\subsection{Experiment 1: Baseline Local-Repair Bias}
\textbf{Setup:} We prompt MiMo-V2.5 with a bug description and source code, without any dependency information. The bug involves changing the signature of \texttt{Task.to\_dict()} in one module---this breaks 8 call sites across 4 files.
\bigskip

\noindent\textbf{Result:} The model identified 3/4 files (75\%), missing \texttt{utils/validators.py}---the most distant dependency from the error traceback. This confirms the local-repair bias.

\subsection{Experiment 2: Visual Diagram (No Text Key)}
\textbf{Setup:} Same bug, but accompanied by a color-coded dependency graph. MiMo reads module names from the diagram using its vision encoder. Qwen2-VL-7B is tested both before and after fine-tuning.
\bigskip

\begin{table}[h]
\centering
\caption{Diagram-only: file coverage on 19 hard test cases}
\begin{tabular}{lccc}
\toprule
Model & Text Only & With Diagram & $\Delta$ \\
\midrule
Base Qwen2-VL-7B  & 4/108 (4\%) & 3/108 (3\%) & --1 \\
Fine-tuned (no CoT) & 22/108 (20\%) & 19/108 (18\%) & --3 \\
\textbf{Fine-tuned (CoT)} & \textbf{14/108 (13\%)} & \textbf{23/108 (21\%)} & \textbf{+9} \\
MiMo-V2.5 (API) & -- & 4/4 (100\%) & -- \\
\bottomrule
\end{tabular}
\end{table}

\textbf{Result:} MiMo's superior vision encoder reads diagram labels perfectly (4/4). Qwen2-VL-7B cannot reliably OCR graph labels, producing hallucinated module names. However, \textbf{CoT fine-tuning reverses the diagram's effect from negative (--3) to positive (+9)}, demonstrating that the model learns to \textit{attempt} diagram-based reasoning even when OCR is imperfect.

\subsection{Experiment 3: Diagram + Minimal Text Key}
\textbf{Setup:} We augment the diagram with a 3-line text key mapping module names to dependency roles: \texttt{"models/task.py (red) -> services/task\_service.py, ... (yellow)"}. This eliminates the OCR requirement while preserving the spatial relationship information from the diagram.
\bigskip

\begin{table}[h]
\centering
\caption{Diagram + Key: toy project generalization test (4 expected files)}
\begin{tabular}{lccc}
\toprule
Model & Text Only & Diagram + Key \\
\midrule
Base Qwen2-VL-7B  & 0/4 & 4/4 \\
Fine-tuned (CoT)  & 0/4 & 4/4 \\
\bottomrule
\end{tabular}
\end{table}

\textbf{Result:} With the text key, both base and fine-tuned models achieve 100\% file coverage on the unseen toy project. The diagram provides spatial intuition; the key provides lexical accuracy. This dual-input architecture---visual graph + compact text legend---is the most robust configuration across model scales.

\subsection{Experiment 4: Legend Distraction}
\textbf{Setup:} We test whether a \textit{full} text legend (listing all 13 modules and all dependencies) helps or hurts. MiMo is tested on diagram-only, text-legend-only, and diagram+legend.
\bigskip

\begin{table}[h]
\centering
\caption{Legend length effect (MiMo-V2.5, toy project)}
\begin{tabular}{lc}
\toprule
Input Mode & Files Found \\
\midrule
Text only (no diagram) & 0/4 \\
Full text legend only & 0/4 \\
\textbf{Diagram only} & \textbf{4/4} \\
Diagram + full legend & 0/4 \\
Diagram + 3-line key & 4/4 \\
\bottomrule
\end{tabular}
\end{table}

\textbf{Result:} A full dependency legend \textit{eliminates} the diagram's benefit, reducing coverage from 4/4 to 0/4. We hypothesize that long text legends displace the model's attention on the visual input, essentially converting a multi-modal task into a pure-text task where the model has insufficient information. The 3-line key preserves both modalities.

\subsection{Experiment 5: CoT Habit Transfer}
\textbf{Setup:} We compare base vs. fine-tuned models on text-only reasoning to test whether the ``describe the diagram'' CoT training transfers to a general structured-analysis habit.
\bigskip

\noindent\textbf{Result:} The fine-tuned model spontaneously produces structured analysis (``Step 1: Identify the changed module... Step 2: Identify dependents...'') even when no diagram is present. This confirms that the CoT training internalizes a \textit{general reasoning pattern}, not merely a diagram-specific template.

% ============================================================
\section{Discussion}

\subsection{Why Visual Graphs Work}
Text is inherently sequential; dependency graphs are inherently parallel. When a model reads ``A depends on B, which calls C, which imports D'' as prose, it must reconstruct the topology through numerous attention hops. A color-coded diagram presents the same information in a single parallel frame. This aligns with how human engineers actually reason about system architecture---through whiteboard sketches, not prose lists.

\subsection{The OCR Bottleneck}
Small vision-language models (7B parameters) have insufficient visual resolution for OCR on graph labels. This emerged as the primary limitation of our approach. However, two practical solutions exist: (1) larger vision models (72B+), and (2) the dual-input architecture with a minimal text key. Option 2 is immediately deployable with any multi-modal LLM.

\subsection{Engineering Cybernetics Connection}
The original motivation for this work drew from \textit{Engineering Cybernetics} (Tsien, 1954): software engineering is fundamentally a problem of controlling complexity through feedback, observability, and system decomposition. Dependency graphs operationalize ``observability''---making the impact surface of a change immediately visible. Training models to consult this representation before acting mirrors the cybernetic principle of closed-loop control: observe $\rightarrow$ analyze $\rightarrow$ act, rather than the open-loop impulse of ``see error $\rightarrow$ fix line.''

\subsection{Limitations}
\begin{enumerate}[nosep]
\item \textbf{Project scale:} Tested on projects up to 100 modules; very-large codebases (1000+ modules) may produce unreadable graphs.
\item \textbf{Language coverage:} Training data is Python-only. The CLI tool supports 5 additional languages but they were not included in fine-tuning.
\item \textbf{Bug diversity:} Our synthetic bugs (parameter/return-type changes) are a subset of real-world bug types.
\end{enumerate}

% ============================================================
\section{Conclusion}

This paper demonstrates that a simple, low-cost intervention---a color-coded dependency graph with a minimal text key---substantially improves LLM performance on cross-module bug-fixing tasks, raising file coverage from 0\% to 100\% in our hardest evaluation setting. Chain-of-Thought fine-tuning on diagram descriptions reverses negative diagram effects (+9 net improvement) and transfers structured-analysis habits to text-only settings. We release \texttt{depgraph}, an open-source multi-language CLI tool that auto-generates dependency graphs and prompt-ready context snippets, enabling immediate adoption in any AI-assisted programming workflow.

Future work includes scaling to very-large codebases (via hierarchical graph decomposition), extending training data to non-Python languages, and testing on larger vision-language models (72B+) to eliminate the OCR bottleneck. A further direction, implied by the CoT findings, is closing the loop: once a model internalizes the ``read diagram $\rightarrow$ describe topology $\rightarrow$ reason'' pipeline, the external diagram becomes a scaffold that can eventually be dropped---the model simulates a visual canvas through structured text. This would realize the original metaphor of a painter who no longer needs to physically step back, because the whole composition is held in working memory.

% ============================================================
\section*{Tool Availability}
\texttt{depgraph} is available at \texttt{https://github.com/yhkkk1234/depgraph} and supports Python, JavaScript, TypeScript, Go, and Rust. Usage:
\begin{verbatim}
  depgraph /path/to/project -o ./output/
  # Generates: dependency_graph.png, dependency_key.txt, prompt_snippet.txt
\end{verbatim}

% ============================================================
% \printbibliography

\end{document}
